\section{C4 — 요인 하나만 흔든 짝 실험}

\claimx{요인 다섯을 하나씩만 흔들어 유보 예측 성능의 $\Delta \pm SE$ 를 냈다. 이 격자 안에서의 순위는 특징 집합(architecture) $>$ 학습 창(curriculum) $>$ 원천 혼합(data mixture) $>$ 도메인 가중(objective) $>$ 능형 $\lambda$(optimizer) 다}{한 수준이 밑판과 두 칸 이상 다르거나, 유보 집합이 수준마다 달라지거나, $\Delta$ 에 짝지은 $SE$ 가 안 붙으면 이 주장은 떨어진다. 격자를 넓히면 스팬이 커지므로 순위는 격자 밖으로 일반화되지 않는다.}

밑판은 HPLT 전량 학습($n=37,222$ 행) · $\lambda=1.0$ · 특징 전량6 · 창 전량 · 균등 가중이고, 유보는 모든 수준에서 같다($n=782$ · 게이트 도메인 합 $728$). 밑판 묶음 유보 스피어만은 $0.37815$ 다.

\begin{table}[h]\centering
\begin{tabular}{llll}
\hline
순위 & 요인 & 스팬 & $\Delta \pm SE$ \\
\hline
1 & R2 특징 집합 (architecture) & 0.098849 & -0.098849 ± 0.041639 \\
2 & R3 학습 창 (curriculum) & 0.071424 & -0.071424 ± 0.021517 \\
3 & R5 원천 혼합 (data mixture) & 0.046575 & -0.046575 ± 0.023094 \\
4 & R4 도메인 가중 (objective) & 0.039736 & -0.039736 ± 0.018558 \\
5 & R1 능형 $\lambda$ (optimizer) & 0.017251 & +0.015264 ± 0.012209 \\
\hline
\end{tabular}
\caption{요인 순위. 스팬은 사전등록된 격자 안에서의 최대--최소 차다.}
\end{table}

문턱(사전등록: $|\Delta|$ 가 두 배 $SE$ 이상이고 카드 문턱 이상)을 둘 다 넘은 수준은 6 / 18 이고, 여섯 전부 $\Delta$ 가 음수다. 즉 이 격자에서 잰 것은 성능을 올리는 요인이 아니라 \emph{망가뜨리면 내려가는} 요인이다.

\section{언급 정밀도의 추정량 정정}

\claimx{등할당 층화 표본에서 층을 무시한 비율은 편향 추정량이고, 제목이 군집이면 독립 시행 가정의 구간은 좁다}{층가중 값이 층 무시 값과 같거나, 제목 군집 붓스트랩 구간이 Clopper--Pearson 구간보다 넓지 않으면 이 주장은 떨어진다.}

층가중 개체 정밀도는 $0.771566$ 이고 층을 무시한 값은 $0.75431$ 다. 라벨한 $696$ 행이 제목 $164$ 개에서 나오므로 독립 시행 가정의 Clopper--Pearson 구간은 좁다: 제목 군집 붓스트랩의 95\% 구간은 [0.559859, 0.872423] 이고 설계효과는 $22.5251$ 다. 유효 삼중쌍은 $27499.4$ [19953.9, 31094] 다.
